\clearpage
\chapter*{第221次报告勘误}
\markboth{第221次报告勘误}{第221次报告勘误}
\addcontentsline{toc}{chapter}{第221次报告勘误}
\sid{e221}
\setcounter{footnote}{3}
\src{302}

\noindent\textbf{6. TDTE IV：Hilbert 概形。}

\begin{center}
[布尔巴基讨论班，第13卷，1960/61年度，第221号，28页。]
\end{center}

\noindent\textbf{第6页，定理2.1，第3行。} --- 把“必要”改为“必要且充分”。

\noindent\textbf{第15页，第13行。} --- 把“$P\neq Q$”改为“$P<Q$”。

\noindent\textbf{第15行。} --- 把“若相反，$P=Q$”改为
“在相反情形下，将有 $P(n)\geq Q(n)$（当 $n$ 足够大时）”。

\noindent\textbf{第18行，公式 (*)。} --- 把“$<$”和“$>$”改为
“$\leq$”和“$\geq$”。

\noindent\textbf{倒数第5行。} --- 把“$r$”改为“$Z$”。

\noindent\textbf{第18页，评注3.9。} --- 应指出，Hartshorne 已研究代数闭域
上的 Hilbert 概形的连通分支。他证明 $\operatorname{Hilb}^{P}_{r}$ 均连通，
并确定了各对 $(r,P)$ 中哪些满足
$\operatorname{Hilb}^{P}_{r}\neq\varnothing$\footnote{HARTSHORNE (R.). --- Connectedness of the Hilbert
scheme (Thesis, Harvard, 1962).}。

\noindent\textbf{第23页，倒数第5行。} --- 把法文“une composante”改为
“d'une composante”（即把“一个分支”改为“一个分支的”）。

\noindent\textbf{第23、24页，评注5.5。} --- 关于 Zappa 的例子，应指出，
Mumford 刚刚构造了 $\mathbf P^3_{\mathbf C}$ 的 Hilbert 概形的一个不可约
分支；其一般点表示次数为14、亏格为24的非奇异曲线，而该分支在其一般点处
非约化。沿所得曲线中的一条作爆破，他还得到一个 $\mathbf C$ 上的三维正则
射影概形；其形式模概形在一般点处非约化，或者等价地，其在 $\mathbf C$ 上
解析几何意义下的局部模簇（参见 Cartan 讨论班，第13卷，1960/61年度）在
每一点处均非约化。

\par
